\section{Verbale della riunione}

\subsection{Risposta al proponente per concordare incontro}

Il proponente di Sanmarco Informatica è stato contattato ed il gruppo ha fissato l'incontro sincrono per il giorno 16 Marzo alle ore 16:00.


\subsection{Verifica dei periodi trascorsi}

Le criticità principali che sono state rilevate:
\begin{itemize}
	\item Un'importante occupazione oraria per quanto riguarda la ristrutturazione dei documenti \textsc{Piano di Progetto}, \textsc{Piano di Qualifica} (con realizzazione di cruscotto web) a seguito delle correzioni rilevate in sede di RR. Anche il Way Of Working è stato migliorato: questo lavoro però ha rallentato e ritardato lo svolgimento delle attività proprie di questi periodi, come la realizzazione del PoC;
	\item Realizzazione del PoC onerosa e iniziata tardi a causa delle modifiche migliorative di cui sopra: nonostante fosse stato scelto l'ultimo giorno disponibile, la realizzazione del PoC si è conclusa troppo a ridosso del colloquio sulla Technology Baseline, rendendo difficoltosa la preparazione alla presentazione stessa;
	\item Poca cura delle attività proprie del ruolo da parte di alcuni componenti: è stato ribadito come la collaborazione di tutti sia necessaria per un buono svolgimento del progetto. In particolare, è importante che ogni componente, conscio del proprio ruolo, si interessi e curi con attenzione le attività che lo riguardano, per evitare che tutti i componenti debbano prendersi carico di tutti i lavori. 
\end{itemize}


// MA DEVO FARE QUESTO??

\subsection{Organizzazione slide per la presentazione RP}

\'E stata strutturata la presentazione per la Revisione di Progettazione nelle sue parti principali, incaricandone i componenti per la realizzazione.

\subsubsection{Introduzione e sistemazione documenti}
\begin{enumerate}
	\item Introduzione modifiche Analisi dei Requisiti per mostrare cambiamenti nella comprensione degli obiettivi di progetto e del nostro specifico contesto: Sofia Chiarello e Giada Zuccolo;
	\item Reimplementazione glossario per semplificarne la consultazione e aggiunta di sezione per la continuous integration dei documenti dalla repository GitHub ad una cartella di Google Drive: Nicolò Greggio;
	\item cruscotto web, sia lato contabilità interna (Google Sheet) sia esposizione dei dati pubblici (Google Sites): Andrea Tessari;
	\item Sistemazione \textsc{Piano di Qualifica} con aggiunta di nuove metriche;
	 
	
\end{enumerate}

